\section{The D0 criterion and the retained orientation data of the Pfaffian--Dirac theorem}
\label{app:obstruction-discharges}

The Pfaffian part of the cross-level identity
\textcircled{\scriptsize 2}\(\leftrightarrow\)\textcircled{\scriptsize 3}
of Chapter~\ref{ch:pfaffian-dirac} is conditional on four obstruction
classes of \S\ref{subsec:ch6-obstructions},
\((D0)\), \((O1)\), \((O1)^+\), \((O2)\), and on the finite geometric
Pfaffian datum \((P_{\mathrm{fin}})\).
The retained D0-HN orientation datum belongs to the
trivial-canonical threefold \(X=K3\times E\); the retained
quotient-orientation datum belongs to \((O1)\); the \((O1)^+\)
transport is relative to that datum; the retained wall-atlas datum
belongs to \((O2)\).
The primitive-recognition datum and the
level-\(\mathsf A\) gravity-line problem remain separate.  The
retained D0 datum, the orientation data, and the finite Pfaffian
sections combine
Joyce--Upmeier orientation theory~\cite{JoyceUpmeier},
Kiem--Li cosection localisation~\cite{KiemLi2013}, the
perverse-sheaf chart cocycle of
Brav--Bussi--Dupont--Joyce--Szendr\H{o}i~\cite{BBDJS2015}, and the
K3-side reduced theory of
Maulik--Pandharipande~\cite{MaulikPandharipande}, with
\(K3\times E\)-specific finite-stabilizer criteria at the Klein-four
edge, the full two-primary \(E[N]\)-edge, half-Hilbert orbit transport,
and the \((R1)\)--\((R2)\) scalar separation.

\subsection{\texorpdfstring{\(K3\times E\)}{K3 x E} datum and criterion theorem}
\label{app:G-setup}

The standing data are: \(S\) a smooth complex projective \(K3\)
surface, \(E\) a smooth elliptic curve, \(X=S\times E\) the
trivial-canonical threefold, \(\sigma_S\in H^0(S,K_S)\) the holomorphic
two-form trivialising \(K_S\), \(\sigma=\mathrm{pr}_S^*\sigma_S\in
H^0(X,p_S^*K_S)\) the inflated K3 two-form defining the
semiregularity cosection, and
\(\mathfrak M^{\mathrm{red},+}_c(X)\) the cosection-reduced moduli
of effective dyonic class \(c\in\Gamma_{\mathrm{eff}}\) on \(X\) in
the chosen stability heart fixed in
Chapter~\ref{ch:k3xe-setup}.  The D0 sector is the
zero-dimensional sector on the threefold \(X\); the Hilbert--Chow
morphism
\(\nu_n:\mathrm{Hilb}^n(K3)\to\mathrm{Sym}^n(K3)\) furnishes only the
K3-fibre scalar test for the reduced theory.  The
Mukai pairing of
Definition~\ref{def:mukai-lattice} is symmetric of signature
\((4,20)\).

\begin{definition}[Retained type-II wall-atlas datum]
\label{def:G-O2-atlas}
The datum \((O2_{\mathrm{atlas}})\) consists, for the three simple
type-II walls \(\delta_1,\delta_2,\delta_3\), of retained wall objects
with full \(\widehat\Gamma_X\)-charge match, semistability on the
wall, reduced self-Ext splitting, rank-one normal Pfaffian block,
invariant Pfaffian unit, quotient orientation, and overlap
compatibility on the cofinal HN system.  Concretely, for each retained
HN height \(R\) it is a system of wall charts
\((U_{\delta,R},u_{\delta,R},\mathscr L_{\delta,R}^{\tan},\upsilon_{\delta,R})\)
with \(s_\delta(u_{\delta,R})=-u_{\delta,R}\),
\[
\operatorname{pf}_{\delta,R}=\upsilon_{\delta,R}u_{\delta,R},
\]
transition isomorphisms on double overlaps, a Cech \(2\)-cocycle
\(c^{\mathrm{wall}}_{\alpha\beta\gamma}\in\{\pm1\}\), and a
\(1\)-cochain \(b^{\mathrm{wall}}_{\alpha\beta}\) satisfying
\(\delta b^{\mathrm{wall}}=c^{\mathrm{wall}}\).  These data are
required to commute with Thom--Sebastiani products, HN transition maps,
the chosen Weyl lifts, and the half-Hilbert orbit chart system carrying
the twelve binary sign coordinates used in the scalar wall count.
\end{definition}

\paragraph{Finite proof tables for \((O2_{\mathrm{atlas}})\).}
At finite HN height \(R\), the retained wall atlas is supplied by the
tables
\[
\mathrm{wall\_objects},\quad
\mathrm{wall\_charts},\quad
\mathrm{overlaps},\quad
\mathrm{orbit\_transport},\quad
\mathrm{half\_hilbert\_orbit},\quad
\mathrm{compatibilities},\quad
\mathrm{scalar\_separation}.
\]
The table \(\mathrm{wall\_objects}\) records full charge match, wall
semistability, reduced self-Ext splitting, the rank-one normal block,
the invariant unit, and the quotient orientation.  The table
\(\mathrm{wall\_charts}\) records normal rank \(1\), divisor order
\(1\), reflection \(u\mapsto -u\), unit invariance, and local sign
\(-1\).  The tables \(\mathrm{overlaps}\) and
\(\mathrm{orbit\_transport}\) record vanishing Cech/cochain,
unit-overlap, orientation-overlap, Pfaffian-transport,
rank-preservation, and sign-preservation defects.  The table
\(\mathrm{half\_hilbert\_orbit}\) records twelve binary coordinates,
orbit cardinality \(2^{12}=4096\), component and boundary coverage,
and zero coverage defect.  The table \(\mathrm{compatibilities}\)
records the Thom--Sebastiani, HN-transition, Weyl-lift,
quotient-orientation, protected-integration, boundary, and component
identities.  The table \(\mathrm{scalar\_separation}\) records the
integer equality \(2^{12}=64^2=4096\) only as a scalar-shadow equality
and excludes the OP scalar branch, the Maass character value, and the
squared theta-leading constant as entries of the wall atlas.  A
single rank-one chart does not supply \((O2_{\mathrm{atlas}})\).
In the present wall-atlas packet these rows are not supplied.  The
finite obstruction ledger
\[
\texttt{certificates/wall\_atlas/k3e\_o2\_atlas/blocked\_obligations.csv}
\]
records the missing wall-object, wall-chart, overlap, Weyl-transport,
half-Hilbert-orbit, compatibility, type-I-exclusion, D0-limit, and
scalar-separation rows.  Its verifier returns
\(\texttt{O2\_OBSTRUCTION\_LEDGER\_VERIFIED}\), with
\(\texttt{o2\_certification=false}\) and
\(\texttt{mathematical\_certification=false}\).  Thus the packet is a
finite absence ledger for \((O2_{\mathrm{atlas}})\), not a construction
of \((O2_{\mathrm{atlas}})\).

\begin{definition}[Retained quotient-orientation datum]
\label{def:G-O1-quot}
The datum \((O1_{\mathrm{quot}})\) consists, on every retained object,
extension, wrapped, mixed, two-step flag, overlap, and HN-transition
stratum, of null-trivialisations of the ordinary reduced orientation
gerbe, the connected \(BE\)-edge class, the finite-stabilizer
Cech--Borel classes, the zero degree-one stabilizer linearizations,
and the Thom--Sebastiani compatibilities of these choices.
\end{definition}

\begin{definition}[Retained finite geometric Pfaffian datum]
\label{def:G-finite-pfaffian-section}
The datum \((P_{\mathrm{fin}})\) is a cofinal HN system carrying the
finite geometric Pfaffian datum of
Definition~\ref{def:finite-geometric-pfaffian-datum}: cosection-reduced
skew perfect complexes, Pfaffian lines, finite Pfaffian sections,
type-II wall divisor orders, automorphic line comparisons, support and
parity rows, normalized leading coefficients, and Pfaffian-line
Mittag--Leffler compatibilities on the cofinal HN system.
In the present certificate packet these rows are not supplied.  The
finite obstruction ledger
\[
\texttt{certificates/pfaffian/k3e\_finite\_pfaffian/blocked\_obligations.csv}
\]
records the missing retained-stratum, skew-complex, Pfaffian-line,
Pfaffian-section, wall-chart, automorphic-comparison, transition, and
scalar-firewall rows.  Its verifier returns
\(\texttt{PFAFFIAN\_OBSTRUCTION\_LEDGER\_VERIFIED}\), with
\(\texttt{pfaffian\_certification=false}\) and
\(\texttt{mathematical\_certification=false}\).  Thus it is a finite
absence certificate for \((P_{\mathrm{fin}})\), not a construction of
\((P_{\mathrm{fin}})\).
\end{definition}

\begin{definition}[Retained Pfaffian--Dirac obstruction datum]
\label{def:G-retained-obstruction-datum}
The retained obstruction datum on \(X=K3\times E\) is the tuple
\[
\mathfrak R_X^{\mathrm{DI}}=
\bigl(
D0_{\mathrm{HN}},\,
O1_{\mathrm{quot}},\,
\mathcal T_{W},\,
O2_{\mathrm{atlas}},\,
P_{\mathrm{fin}},\,
Z_{\mathrm{HB}}
\bigr).
\]
Here \(D0_{\mathrm{HN}}\) is the retained D0-HN orientation datum of
Definition~\ref{def:G-D0-HN-system}; \(O1_{\mathrm{quot}}\) is the
quotient-orientation datum of Definition~\ref{def:G-O1-quot};
\(\mathcal T_W\) is the torsor of Weyl determinant-line lifts
\(\{\tau_i\}_{i=1}^3\) of
Definition~\ref{def:O1-plus-weyl-transport}, together with a chosen
\(1\)-cochain killing the projective cocycle \(c_o\), the equations
\(\tau_i^2=1\), and the vanishing torsor defects
\(\omega_{i,\mathcal C}=0\) on every retained orbit;
\(O2_{\mathrm{atlas}}\) is the retained type-II wall atlas of
Definition~\ref{def:G-O2-atlas}; \(P_{\mathrm{fin}}\) is the retained
finite geometric Pfaffian datum of
Definition~\ref{def:G-finite-pfaffian-section}; and \(Z_{\mathrm{HB}}\)
is the level-\(\mathsf Z\) Hall--Borcherds trace comparison datum of
Definition~\ref{def:G-vol2-hall-borcherds}.

The projections of \(\mathfrak R_X^{\mathrm{DI}}\) are part of the
definition.  The Pfaffian and orientation statements use the first five
components.  The level-\(\mathsf Z\) protected trace uses \(Z_{\mathrm{HB}}\).
The primitive-recognition theorem uses neither \(Z_{\mathrm{HB}}\) nor
the mirror discriminant; it requires the finite compact-source
recognition datum \((S1)\)--\((S10)\).  A scalar product, signed
denominator table, protected trace, or mirror period identity is not a
component of \(\mathfrak R_X^{\mathrm{DI}}\).
\end{definition}

\begin{theorem}[Relative Pfaffian-orientation criterion on \(K3\times E\)]
\label{thm:G-four-obstruction-criterion}
On the rank-one and finite-stabiliser strata of the moduli stack
\(\mathfrak M^{\mathrm{red},+}_c(X)\), \(X=K3\times E\), in the
stability heart fixed in Chapter~\ref{ch:k3xe-setup}, let
\(\mathfrak R_X^{\mathrm{DI}}\) be the retained obstruction datum of
Definition~\ref{def:G-retained-obstruction-datum}.  The obstruction
class \((D0)\) is represented by its \(D0_{\mathrm{HN}}\)-component.
The class \((O1)\) is represented by its \(O1_{\mathrm{quot}}\)-component,
\((O1)^+\) by its Weyl-lift torsor \(\mathcal T_W\), and \((O2)\) by
its \(O2_{\mathrm{atlas}}\)-component.  If the Cech, Borel, torsor,
and wall-atlas equations in these components vanish with the stated
null-trivialisations, then the Pfaffian--Dirac theorem
\ref{thm:ch6-pfaffian-dirac} and the orientation-character match
\ref{thm:ch6-orientation-character} hold on \(K3\times E\) relative to
\(\mathfrak R_X^{\mathrm{DI}}\) through its first five components.  The primitive
recognition theorem
\ref{thm:ch6-primitive-recognition} remains conditional on the finite
compact-source recognition datum \((S1)\)--\((S10)\).
\end{theorem}

The four clauses are separated below according to their retained data.
The level-\(\mathsf Z\) Hall--Borcherds trace comparison is distinct
from the level-\(\mathsf A\) gravity-line residual.

\subsection{(D0) D0-degeneration limit of the cosection-reduced d-critical orientation theory}
\label{app:G-D0-discharge}

The first obstruction concerns the inverse-limit assembly of the
Joyce--Upmeier orientation cocycle along the HN system
\(\mathcal R_{\mathrm{HN}}\) when the chosen stability heart degenerates to
the \(D0\)-supported subcategory of zero-dimensional sheaves on
\(X=K3\times E\).  The \(D0\)-charge is the threefold length \(n\).
Projection to the K3 factor gives the reduced K3 scalar test; it does
not replace the threefold D0 moduli.  The D0-degeneration limit must
commute with cosection-reduction and with the orientation extension.

\begin{definition}[Retained D0-HN orientation datum]
\label{def:G-D0-HN-system}
\textnormal{[definitional]}
For each \(R\in\mathcal R_{\mathrm{HN}}\), let
\(\mathfrak M_R\) be the retained finite-type HN substack of objects of
height \(\le R\) in the chosen stability heart, and let
\(\mathfrak M_R^{D0}\subset\mathfrak M_R\) be the retained
zero-dimensional threefold sector, with \(D0\)-charge equal to length
on \(X=K3\times E\).  The datum \((D0_{\mathrm{HN}})\) consists of:
\begin{enumerate}
\item[\textup{(D0.1)}] A flat one-parameter degeneration
\(\mathfrak M_R\to \mathbb A^1\) whose special fibre contains
\(\mathfrak M_R^{D0}\), together with compatible object, extension,
mixed, wrapped, and two-step flag substacks over \(\mathbb A^1\).
\item[\textup{(D0.2)}] Proper or closed HN transition morphisms
\[
\rho^{\mathrm{stk}}_{RR'}:\mathfrak M_R|_{\le R'}
\longrightarrow \mathfrak M_{R'}
\qquad (R'\le R)
\]
on the object, extension, and flag stacks, satisfying the composition
law and preserving the \(D0\)-special fibre.
\item[\textup{(D0.3)}] Quasi-smooth derived enhancements with
PTVV \((-1)\)-shifted symplectic forms, classical d-critical
truncations, and perfect obstruction theories on every retained
stratum of the degeneration family.
\item[\textup{(D0.4)}] The K3 semiregularity cosection
\[
\operatorname{CS}_{\mathcal C,R}:
h^1\!\left((E^{\bullet}_{\mathcal C,R})^\vee\right)\longrightarrow
\mathcal O_{\mathfrak M_R}
\]
obtained by contracting the obstruction theory with
\(\operatorname{pr}_S^*\sigma_S\), surjective on every retained reduced
stratum, compatible with the \(D0\)-special fibre, and additive on
extension and flag stacks:
\[
q^*\operatorname{CS}_{C,R}
=p_1^*\operatorname{CS}_{A,R}+p_2^*\operatorname{CS}_{B,R}.
\]
\item[\textup{(D0.5)}] Cosection-localized vanishing-cycle complexes
\(\Phi^{\mathrm{red}}_R\), with specialization isomorphisms from the
generic fibre to the \(D0\)-special fibre and transition isomorphisms
under \(\rho^{\mathrm{stk}}_{RR'}\).
\item[\textup{(D0.6)}] Joyce--Upmeier orientation lines
\(o_R\) on the reduced d-critical charts, with square-root
isomorphisms for the reduced determinant complex, Thom--Sebastiani
multiplicativity on retained extensions, and transition isomorphisms
carrying \(o_R\) to \(o_{R'}\).
\item[\textup{(D0.7)}] Pfaffian-line and protected-integration
compatibility: the Pfaffian lines, finite Pfaffian sections,
quotient-after-correspondence maps, compact-support six-functor
operations, and protected integration maps extend across the
\(D0\)-special fibre and commute with \(\rho^{\mathrm{stk}}_{RR'}\).
\item[\textup{(D0.8)}] Strict Mittag--Leffler exactness for the inverse
systems of retained moduli cohomology, reduced vanishing cycles,
orientation gerbes, Pfaffian lines, compact-support operations, Hall
product/coproduct maps, and the corresponding entries of the finite-stage morphism
\(\rho_{RR'}\) of
Definition~\ref{def:finite-stage-dirac-igusa-morphism}.
\end{enumerate}
The D0-degeneration limit is the pro-object obtained from
\textup{(D0.1)}--\textup{(D0.8)}.  The scalar Hilbert-scheme
specialization on \(K3\) is only a test of the reduced scalar theory;
it is not an entry of \((D0_{\mathrm{HN}})\) and does not construct the
transition morphisms above.

After finite presentations of the HN stacks and their correspondence
atlases, \((D0_{\mathrm{HN}})\) is recorded by the D0 proof tables
\[
\mathrm{degeneration\_families},\quad
\mathrm{retained\_substacks},\quad
\mathrm{hn\_transitions},\quad
\mathrm{derived\_enhancements},\quad
\mathrm{semiregularity\_cosections},\quad
\mathrm{vanishing\_cycles},\quad
\mathrm{orientation\_specializations},\quad
\mathrm{pfaffian\_integration},\quad
\mathrm{ml\_exactness},\quad
\mathrm{scalar\_firewall}.
\]
These rows are finite data: flatness defects, composition defects,
obstruction-theory defects, cosection cokernels, additivity defects,
specialization cones, transition defects, and \(R^1\varprojlim\)-ranks
must vanish.  The retained-substack rows must include object,
extension, mixed, wrapped, and two-step flag substacks.  The
Mittag--Leffler rows must include retained moduli cohomology, vanishing
cycles, orientation gerbes, Pfaffian lines, compact-support operations,
Hall product, Hall coproduct, and finite-stage morphisms.  The scalar firewall excludes
Hilbert-scheme scalar specialization, the K3 scalar test, and scalar
trace as entries of \((D0_{\mathrm{HN}})\).
In the present certificate packet the D0-HN rows are not supplied.
The finite obstruction ledger
\[
\texttt{certificates/d0/k3e\_d0\_hn/blocked\_obligations.csv}
\]
records the missing flat-family, D0-sector, retained-substack,
HN-transition, derived-enhancement, semiregularity-cosection,
vanishing-cycle, orientation, Pfaffian, protected-integration,
Mittag--Leffler, and scalar-firewall rows.  Its verifier returns
\(\texttt{D0\_OBSTRUCTION\_LEDGER\_VERIFIED}\), with
\(\texttt{d0\_certification=false}\) and
\(\texttt{mathematical\_certification=false}\).  Thus it is a finite
absence certificate for \((D0_{\mathrm{HN}})\), not a proof of
\((D0_{\mathrm{HN}})\).
\end{definition}

\begin{theorem}[(D0) criterion]
\label{thm:G-D0}
Granted \((D0_{\mathrm{HN}})\), the cosection-reduced d-critical
orientation theory extends compatibly along the HN inclusions
\(\iota_R\) and has a D0-degeneration pro-limit on
\(\mathfrak M^{\mathrm{red}}_\infty\).
\end{theorem}

\begin{proof}
PTVV supplies the ambient \((-1)\)-shifted symplectic source on the
trivial-canonical threefold moduli \cite[Theorem~0.3]{PTVV2013}; BBDJS
constructs vanishing cycles after an oriented d-critical atlas is
present \cite[Theorem~6.9]{BBDJS2015}; Kiem--Li gives the cosection
localisation formalism \cite{KiemLi2013}; Joyce--Upmeier gives the
orientation-line and multiplicativity technology
\cite{JoyceUpmeier,JoyceTanakaUpmeier}.  These theorems do not by
themselves produce the retained HN-compatible reduced orientation
system.  Clauses \textup{(D0.1)}--\textup{(D0.3)} supply the
degeneration family and finite derived/d-critical atlases; clause
\textup{(D0.4)} supplies the surjective additive K3 semiregularity
cosection; clauses \textup{(D0.5)}--\textup{(D0.7)} supply the
vanishing-cycle, orientation, Pfaffian, Hall, quotient, and
protected-integration specializations; and clause \textup{(D0.8)}
supplies the strict Mittag--Leffler exactness needed to pass to the
pro-limit.
The pro-limit therefore inherits the reduced vanishing-cycle complexes,
orientation lines, Thom--Sebastiani transports, compact-support
operations, and Pfaffian lines.  Thus the retained HN datum gives
the \((D0)\) vanishing condition on the pro-limit.
\end{proof}

\begin{remark}[The role of the K3 elliptic genus]
\label{rem:G-D0-elliptic-genus}
The K3 elliptic-genus formula
\(\chi_y(\mathrm{Hilb}^n(K3))\) \cite[\S 4.1]{MaulikPandharipande}
tests the scalar D0 specialization.  It is not a substitute for
Mittag--Leffler exactness of the orientation gerbes, vanishing-cycle
complexes, compact-support functors, or Pfaffian lines.
\end{remark}

\subsection{(O1) strong reduced orientation on the K3\(\times\)E-quotient self-Ext frame bundle}
\label{app:G-O1-discharge}

The second obstruction is the strong reduced orientation: the
multiplicative orientation cocycle of \((D0)\) must be equivariant
under the free \(E\)-translation quotient and under the
finite-stabilizer strata \(BE[N]\) on which the
\(E\)-action degenerates.

\begin{definition}[Strong reduced orientation on the K3\(\times E\)-quotient frame]
\label{def:G-strong-reduced-orientation}
\textnormal{[definitional]}
The \emph{strong reduced orientation} on \(\mathfrak M^{\mathrm{red},+}_c(X)\)
is a multiplicative orientation cocycle
\(o^{\mathrm{strong}}_c\) on the principal \(\mathrm{Spin}^c\)-bundle
\(\mathcal F^{\mathrm{red}}_c\) of the reduced derived self-\(\mathrm{Ext}\)
frame bundle on \(\mathfrak M^{\mathrm{red},+}_c(X)\), satisfying:
\begin{enumerate}
\item[(O1.a)] \emph{Free \(E\)-translation equivariance.}  On the
generic \(E\)-translation orbit of codimension \(0\),
\(o^{\mathrm{strong}}_c\) descends through the free quotient
\(\mathfrak M^{\mathrm{red},+}_c(X)/E\).
\item[(O1.b)] \emph{Klein-four edge and linearization condition.}
Across the \(BE[2]\)-strata, the Cech--Borel class reduces to the
edge class of Lemma~\ref{lem:G-Klein-four}, that class vanishes, and
the residual degree-one \(E[2]\)-linearization character is zero.
\item[(O1.c)] \emph{Full two-primary finite-stabilizer condition.}
Across the \(BE[N]\)-strata for \(N\ge 3\), the two-primary edge class
and the residual stabilizer character satisfy the criterion of
Lemma~\ref{lem:G-two-primary-stabilizer}.
\end{enumerate}
\end{definition}

\begin{lemma}[Klein-four edge criterion on rank-\(2\) \(E\)-strata]
\label{lem:G-Klein-four}
Assume that the finite-stabilizer Cech--Borel class on a retained
\(BE[2]\)-stratum has no ordinary, mixed, or differential term and
therefore has an edge representative in \(H^2(BE[2],\F_2)\).  For
\(E[2]\simeq(\Z/2)^2\),
\[
H^*(BE[2],\F_2)=\F_2[x_1,x_2],\qquad |x_i|=1,
\]
and the edge class has the form
\[
\beta^{E,2}=b_{20}x_1^2+b_{11}x_1x_2+b_{02}x_2^2.
\]
If \(r_i\) denotes its restriction to the three non-zero cyclic
subgroups of \(E[2]\), then
\[
b_{20}=r_1,\qquad b_{02}=r_2,\qquad
b_{11}=r_1+r_2+r_3.
\]
Hence \(\beta^{E,2}=0\) is equivalent to \(r_1=r_2=r_3=0\) for this
edge class.  After this degree-two class is killed, the remaining
linearization character is
\[
\lambda^{E,2}=\lambda_1x_1+\lambda_2x_2\in H^1(BE[2],\F_2),
\]
and its restrictions satisfy
\[
\rho_1=\lambda_1,\qquad \rho_2=\lambda_2,\qquad
\rho_3=\lambda_1+\lambda_2.
\]
The vanishing of \(r_i\) is therefore not the vanishing of
\(\rho_i\); quotient orientation requires both sets of equations.
\end{lemma}

\begin{proof}
The cohomology ring is the polynomial ring displayed above.  Restricting
to \(\langle e_1\rangle\), \(\langle e_2\rangle\), and
\(\langle e_1+e_2\rangle\) sends the quadratic form respectively to
\(b_{20}t^2\), \(b_{02}t^2\), and
\((b_{20}+b_{11}+b_{02})t^2\).  Over \(\F_2\) this gives the displayed
invertible linear change of coordinates.  The last assertion is the
standard statement that, after a degree-two gerbe is trivialised,
linearizations form an \(H^1(BE[2],\F_2)\)-torsor.
\end{proof}

\begin{lemma}[Full two-primary finite-stabilizer criterion]
\label{lem:G-two-primary-stabilizer}
For a complex elliptic curve,
\[
E[N]\simeq(\Z/N)^2.
\]
If \(N\) is odd, the mod-\(2\) finite-stabilizer orientation class
restricts trivially by transfer.  If \(2^a\parallel N\), then the
mod-\(2\) calculation factors through
\[
E[N]_{(2)}\simeq(\Z/2^a)^2.
\]
For \(a=1\) one obtains Lemma~\ref{lem:G-Klein-four}.  For \(a\ge2\),
\[
H^*\bigl(B(\Z/2^a)^2,\F_2\bigr)
=\Lambda(x_1,x_2)\otimes \F_2[y_1,y_2],
\qquad |x_i|=1,\quad |y_i|=2,
\]
and every degree-two edge class has the form
\[
\beta^{E,N}=A_1^{(N)}y_1+A_{12}^{(N)}x_1x_2+A_2^{(N)}y_2.
\]
The residual stabilizer character is
\[
\lambda^{E,N}=\lambda_1^{(N)}x_1+\lambda_2^{(N)}x_2.
\]
Thus quotient orientation at the \(N\)-stratum requires
\[
A_1^{(N)}=A_{12}^{(N)}=A_2^{(N)}=
\lambda_1^{(N)}=\lambda_2^{(N)}=0.
\]
The coefficient \(A_{12}^{(N)}\) is not detected by cyclic order-two
restrictions.
\end{lemma}

\begin{proof}
The group of \(N\)-torsion points of a complex elliptic curve is
\((\Z/N)^2\).  For \(N\) odd the group order is invertible in
\(\F_2\), hence \(H^{>0}(BE[N],\F_2)=0\) by transfer.  The
even-primary reduction is the primary decomposition of \(E[N]\).  The
cohomology of
\(B(\Z/2^a)\) is \(\Lambda(x)\otimes\F_2[y]\), with
\(|x|=1\) and \(|y|=2\); the K\"unneth formula gives the displayed
degree-two basis for the square.  A cyclic order-two subgroup pulls
back \(x_1x_2\) to zero, so it cannot see the mixed coefficient.
\end{proof}

\begin{remark}[What the finite-stabilizer criteria prove]
\label{rem:G-N-2-vs-N-3}
The two preceding lemmas are detection statements inside the retained
datum \((O1_{\mathrm{quot}})\).  They do not construct the Cech--Borel
null-trivialisations, do not kill mixed Borel terms, and do not remove
the residual \(H^1\)-linearization characters.  They say which
coordinates must vanish after the finite-stabilizer problem has been
reduced to the classifying-space edge.
The finite obstruction ledger
\(\texttt{certificates/orientation/k3e\_reduced\_orientation/blocked\_obligations.csv}\)
records the missing orientation rows before this criterion can be used
as an unconditional theorem: square-root, Thom--Sebastiani, quotient
Borel, finite-stabilizer, Weyl-lift, Coxeter, transition, and
scalar-firewall rows.  Its verification status is
\(\texttt{ORIENTATION\_OBSTRUCTION\_LEDGER\_VERIFIED}\), not
orientation certification.
\end{remark}

\begin{theorem}[(O1) quotient-orientation criterion]
\label{thm:G-O1}
Granted the retained quotient-orientation datum
\((O1_{\mathrm{quot}})\) of Definition~\ref{def:G-O1-quot}, the
reduced orientation of Theorem~\ref{thm:G-D0} descends to the strong
reduced orientation \(o^{\mathrm{strong}}_c\) of
Definition~\ref{def:G-strong-reduced-orientation} on
\(\mathfrak M^{\mathrm{red},+}_c(X)\).
\end{theorem}

\begin{proof}
The reduced orientation \(o^{\mathrm{red}}_\infty\) of
Theorem~\ref{thm:G-D0} and the datum \((O1_{\mathrm{quot}})\) determine
\((O1.a)\), \((O1.b)\), and \((O1.c)\).

\smallskip\emph{\((O1.a)\) Free \(E\)-equivariance.}  On the free
translation locus the connected \(BE\)-edge class
\(\alpha^{E,\mathrm{free}}\in H^2(BE,\F_2)\) is one of the
null-trivialised classes in \((O1_{\mathrm{quot}})\).  The chosen
trivialisation, together with Joyce--Upmeier multiplicativity
\cite[\S4]{JoyceUpmeier}, gives descent through the open quotient.

\smallskip\emph{\((O1.b)\) Klein-four extension.}
Lemma~\ref{lem:G-Klein-four} identifies the \(BE[2]\)-edge equations.
The datum \((O1_{\mathrm{quot}})\) supplies their vanishing and the
zero \(H^1(BE[2],\F_2)\)-linearization character.

\smallskip\emph{\((O1.c)\) Higher-rank extension.}
Lemma~\ref{lem:G-two-primary-stabilizer} identifies the full two-primary
finite-stabilizer equations for \(BE[N]\), \(N\ge3\).  The same datum
supplies their vanishing and the compatible zero stabilizer
characters on overlaps, Hall correspondences, flags, and HN
transitions.

The combined cocycle is multiplicative and Pin\(^c\)-equivariant by
the multiplicativity of the supplied Thom--Sebastiani transports; this
is \(o^{\mathrm{strong}}_c\).  By the multiplicativity, every product
\(F_1\to E\to F_2\) of effective dyonic objects has
\(o^{\mathrm{strong}}_c(F_1\to E\to F_2)
=o^{\mathrm{strong}}_c(F_1)\otimes o^{\mathrm{strong}}_c(F_2)\).
This is the strong reduced orientation.
\end{proof}

\subsection{(O1)\(^+\) Weyl-equivariant transport along
\texorpdfstring{\(W^{(2)}(\Lambda^{2,1}_{II})\)}{type-II Weyl group}}
\label{app:G-O1-plus-discharge}

The third obstruction is the Weyl-equivariance of the strong reduced
orientation along the type-II Weyl group
\(W^{(2)}(\Lambda^{2,1}_{II})\).  The Weyl group is the
\(3\)-generator Coxeter group generated by the three simple type-II
reflections \(s_{\delta_1}\), \(s_{\delta_2}\), \(s_{\delta_3}\).  Since
\((\delta_i,\delta_j)=-2\) for \(i\ne j\), the Coxeter exponents are
\(m_{ij}=\infty\); there are no braid relations among distinct
generators.  Hence
\((W^{(2)})_{\mathrm{ab}}\simeq(\Z/2)^{\oplus 3}\).

\begin{theorem}[(O1)\(^+\) Weyl transport criterion and type-II character comparison]
\label{thm:G-O1-plus}
The strong reduced orientation
\(o^{\mathrm{strong}}_c\) of Theorem~\ref{thm:G-O1}, together with the
chosen lifts of Definition~\ref{def:O1-plus-weyl-transport} with
vanishing projective cocycle and zero quotient-torsor defects, is
\(W^{(2)}(\Lambda^{2,1}_{II})\)-equivariant.  If, in addition, the
rank-one type-II wall signs of \((O2_{\mathrm{atlas}})\) are fixed on
the three simple walls, then its character on
\(W^{(2)}\) is
\[
 \epsilon_o\;=\;\nu_{\Delta_5}\quad\text{on }W^{(2)}(\Lambda^{2,1}_{II}),
\]
identical to the Maass multiplier system of the Igusa cusp form.
\end{theorem}

\begin{proof}
The proof compares the equivariant orientation line with the Maass
character on the type-II Coxeter generators.

\smallskip\emph{Equivariant determinant line.}
The vanishing of the projective cocycle \(c_o\) makes the chosen
determinant-line lifts into an honest action of
\(W^{(2)}(\Lambda^{2,1}_{II})\) on \(o^{\mathrm{strong}}_c\).  The
vanishing of the torsor defects \(\omega_{i,\mathcal C}\) says that the
quotient-orientation cocycle is transported with the line, not only
after forgetting the finite-stabilizer data.  This proves the
\((O1)^+\) equivariance assertion.

\smallskip\emph{Generator values from \((O2)\).}
Under \((O2_{\mathrm{atlas}})\), Proposition~\ref{prop:appE-local-pfaffian-sign}
gives the local Pfaffian sign at each simple type-II wall:
\[
 \epsilon_o^{\mathrm{loc}}(s_{\delta_i})\;=\;
 (-1)^{N_{\delta_i}^{\mathrm{Pf}}}\;=\;-1,
 \qquad i=1,2,3,
\]
with \(N_{\delta_i}^{\mathrm{Pf}}=1\) the rank-one Pfaffian rank
exponent.  By Lemma~\ref{lem:appE-maass-on-type-ii}, the Maass
character on these generators is
\(\nu_{\Delta_5}(s_{\delta_i})=-1\) for \(i=1,2,3\).
Hence \(\epsilon_o(s_{\delta_i})=\nu_{\Delta_5}(s_{\delta_i})=-1\)
for the three generators.

\smallskip\emph{The abelianisation \((\Z/2)^{\oplus 3}\).}
The type-II Weyl group has the Coxeter presentation
\[
W^{(2)}(\Lambda^{2,1}_{II})
=\langle s_1,s_2,s_3\mid s_1^2=s_2^2=s_3^2=1\rangle,
\]
with no relation on \(s_is_j\) for \(i\ne j\).
This is the Coxeter presentation attached to the symmetric matrix
\((\delta_i,\delta_j)\), because \((\delta_i,\delta_j)=-2\) gives
\(m_{ij}=\infty\).  Abelianising only imposes commutativity and
\(s_i^2=1\), hence
\[
(W^{(2)})_{\mathrm{ab}}\simeq(\Z/2)^{\oplus 3}.
\]

\smallskip\emph{Extension to the full Weyl group.}
A character of \(W^{(2)}\) is determined by its values on the
abelianisation, hence by its values on the three generators.  The
extension principle is standard: given
\(\chi:\{s_1,s_2,s_3\}\to\{\pm 1\}\), define
\(\chi(w)=\prod_{j=1}^k\chi(s_{i_j})\) for any reduced word
\(w=s_{i_1}\cdots s_{i_k}\).  The only defining relations are
\(s_i^2=1\); hence the formula is well-defined.

\smallskip\emph{Character comparison.}
The group \(H^1(W^{(2)},\{\pm 1\})\) is
\(\operatorname{Hom}(W^{(2)},\{\pm1\})\simeq(\Z/2)^{\oplus3}\), not the
obstruction group.  The obstruction is represented by the
projective cocycle \(c_o\) and torsor defects \(\omega_{i,\mathcal C}\)
of Definition~\ref{def:O1-plus-weyl-transport}.  When these classes
vanish and the three rank-one \((O2)\) signs are fixed,
both \(\epsilon_o\) and \(\nu_{\Delta_5}\) are
\(\{\pm 1\}\)-valued characters of \(W^{(2)}\) with the same values on
the three generators by the local \((O2)\) computation; hence they are equal:
\(\epsilon_o=\nu_{\Delta_5}\) on \(W^{(2)}(\Lambda^{2,1}_{II})\).
\end{proof}

\begin{remark}[Why the type-II Coxeter group reduces to its abelianisation]
\label{rem:G-O1-plus-reduction}
The reduction to the abelianisation is forced by the order-two structure of the
\(\{\pm 1\}\)-character group: any group homomorphism into a group of
exponent two factors through the abelianisation, since exponent-two
characters automatically annihilate the commutator subgroup.  There is
no order-three braid relation in the type-II Coxeter group; the
off-diagonal value \(-2\) is the infinite Coxeter exponent.  Thus no
order-three obstruction class is present in this comparison.
\end{remark}

\subsection{(O2) Local Pfaffian wall normal form and the \((R1)\)--\((R2)\) scalar separation}
\label{app:G-O2-discharge}

The fourth obstruction is the type-II wall normal form: the local
Pfaffian section on each type-II wall is the rank-one odd
crossing of Appendix~E, and the \(2^{12}=4096\) sign-sum on the
half-Hilbert orbit remains separate from the OP normalization until
the scalar shadow is taken.

\begin{theorem}[(O2) wall-atlas consequence: type-II wall normal form]
\label{thm:G-O2-normal-form}
Assume \((O2_{\mathrm{atlas}})\) and the Weyl-transport datum of
Theorem~\ref{thm:G-O1-plus}.  On each type-II wall
\(\delta\in W^{(2)}(\Lambda^{2,1}_{II})\cdot
\{\delta_1,\delta_2,\delta_3\}\) of
\(\mathfrak M^{\mathrm{red},+}_c(X)\), the local Pfaffian section
\(\operatorname{pf}_\delta\) is the rank-one odd crossing
\([\R\xrightarrow{u}\R]\) with Pfaffian rank exponent
\(N_\delta^{\mathrm{Pf}}=1\) and local sign
\(\epsilon_o^{\mathrm{loc}}(s_\delta)=-1\).
\end{theorem}

\begin{proof}
The proof reduces higher-rank type-II walls to the rank-one normal
form of Appendix~E by Weyl-orbit transport.

\smallskip\emph{Rank-one walls.}
For \(\delta\in\{\delta_1,\delta_2,\delta_3\}\), the rank-one normal
form is Proposition~\ref{prop:appE-local-pfaffian-sign}, supplying
\(N_\delta^{\mathrm{Pf}}=1\) and
\(\epsilon_o^{\mathrm{loc}}(s_\delta)=-1\).

\smallskip\emph{Orbit-stabiliser theorem.}
For higher-rank walls
\(\delta=w\cdot\delta_i\) with \(w\in W^{(2)}(\Lambda^{2,1}_{II})\)
and \(i\in\{1,2,3\}\), the local geometric chart
\(U_\delta=T_\delta\times(-\varepsilon,\varepsilon)\) is identified in
the retained wall atlas with the \(w\)-transport of the rank-one chart
\(U_{\delta_i}=T_{\delta_i}\times(-\varepsilon,\varepsilon)\).  By
Theorem~\ref{thm:G-O1-plus} (Weyl-equivariance of the strong reduced
orientation), the local Pfaffian section transports as
\[
 \operatorname{pf}_\delta\;=\;w_*\,\operatorname{pf}_{\delta_i}
 \;\cdot\;\epsilon_o(w),
\]
with \(\epsilon_o(w)\in\{\pm 1\}\) the Weyl character value.

\smallskip\emph{Rank exponent invariance.}
The Pfaffian rank exponent
\(N_\delta^{\mathrm{Pf}}=1\) is preserved under the supplied
Weyl-orbit chart identifications, since it is the rank of the normal
factor in the reduced d-critical chart.

\smallskip\emph{Local sign.}
The local sign \(\epsilon_o^{\mathrm{loc}}(s_\delta)=
\epsilon_o^{\mathrm{loc}}(w s_{\delta_i}w^{-1})\).  Since the
orientation sign is a character after Theorem~\ref{thm:G-O1-plus},
\[
\epsilon_o(w s_{\delta_i}w^{-1})
=\epsilon_o(w)\epsilon_o(s_{\delta_i})\epsilon_o(w)^{-1}
=\epsilon_o(s_{\delta_i})=-1.
\]
Hence
\(\epsilon_o^{\mathrm{loc}}(s_\delta)=-1\) for every Weyl-translate.
\end{proof}

\begin{theorem}[\((R1)\)--\((R2)\) scalar separation: \(4096=2^{12}=64^2\)]
\label{thm:G-R1-R2}
Two independent normalizations have the same integer value:
\begin{description}
\item[\textnormal{(R1)}] (operator-level)
\(N_{\mathrm{wall}}=2^{12}\), the count of sign assignments in the
half-Hilbert wall atlas of \textup{(O2)};
\item[\textnormal{(R2)}] (automorphic-level)
\(N_{\mathrm{OP}}=64^2=c_\Delta(1,1,1)^2\), the squared theta-leading
constant of the type-II Borcherds product.
\end{description}
The equality \(N_{\mathrm{wall}}=N_{\mathrm{OP}}=4096\) is an equality
of integers after passing to a scalar shadow.  It is not an
identification of the orientation character with the OP scalar
normalization.
\end{theorem}

\begin{proof}
The (O2) atlas supplies \(12\) binary wall signs, hence
\[
N_{\mathrm{wall}}=2^{12}=4096.
\]
The leading Fourier coefficient of \(\Delta_5\) at the type-II
vacuum is
\[
 c_\Delta(1,1,1)\;=\;[q^{1/2}r^{1/2}s^{1/2}]\,\Delta_5
 \;=\;64
\]
by Lemma~\ref{lem:bkm-vacuum-coefficient}.  Squaring,
\[
N_{\mathrm{OP}}=c_\Delta(1,1,1)^2=64^2=4096.
\]
Thus
\[
N_{\mathrm{wall}}=N_{\mathrm{OP}}.
\]
No step identifies the wall orientation line with the OP scalar branch.
\end{proof}

\begin{remark}[\((R1)\) and \((R2)\) live on different levels]
\label{rem:G-R1-R2}
The wall count belongs to the Pfaffian orientation atlas.  The factor
\(64^2\) belongs to the automorphic normalization
\(D_5=64^{-1}\Delta_5\) and to the OP scalar convention.  Their common
integer is used only after orientation has been forgotten by squaring.
\end{remark}

\begin{theorem}[(O2) wall-atlas transport on the half-Hilbert orbit]
\label{thm:G-O2-orbit-transport}
Assume \((O2_{\mathrm{atlas}})\) and the Weyl-transport datum of
Theorem~\ref{thm:G-O1-plus}.
The orientation cocycle on the half-Hilbert orbit on
\(\mathfrak M^{\mathrm{red},+}_c(X)\) extends from the rank-one local
data of Theorem~\ref{thm:G-O2-normal-form} to a globally consistent
\((\Z/2)\)-valued system on the entire half-Hilbert orbit, with the
total sign count \(2^{12}\) realised by the twelve binary wall
coordinates of the \((O2_{\mathrm{atlas}})\) datum.  Equivalently, the type-II
Pfaffian-wall normal form is reduced to \((O2_{\mathrm{atlas}})\).
\end{theorem}

\begin{proof}
Theorem~\ref{thm:G-O2-normal-form} supplies the rank-one normal form
on every type-II wall relative to the wall atlas.  Theorem~\ref{thm:G-O1-plus}
supplies the \(W^{(2)}\)-equivariance.  Theorem~\ref{thm:G-R1-R2} identifies the
integer \(2^{12}=4096\) with the squared theta-leading normalization,
after passing to the scalar shadow.
The half-Hilbert orbit transport is the following Weyl-orbit
calculation.  A base point on the orbit determines every other orbit
point by Weyl transport, and the orientation sign is the product of
generator signs along the chosen path.  The result is independent of
path by the chosen Weyl lifts and zero torsor defects of
Theorem~\ref{thm:G-O1-plus}, and the total count is the
\(2^{12}=4096\) count supplied by the wall
atlas.  This is the \((O2)\) transport theorem relative to
\((O2_{\mathrm{atlas}})\).
\end{proof}

\subsection{Vol II Hall--Borcherds trace comparison at level \texorpdfstring{$\mathsf Z$}{Z}: chain-level chiral-Hochschild trace}
\label{app:G-vol2-residual}

The retained D0-HN datum, retained quotient-orientation datum, chosen
Weyl lifts, and retained type-II wall atlas give the relative
Pfaffian and orientation hypotheses of the Pfaffian--Dirac theorem on
\(K3\times E\).  A
separate trace comparison relates the level-\(n\) protected scalar shadow
\(Z_{\mathrm{BPS}}^{K3\times E}=(\Phi_{10}^{\mathrm{un}})^{-1}=
\Delta_5^{-2}\) of Chapter~\ref{ch:zbps-realization} to a chain-level
chiral-bulk trace on the derived chiral centre
\(Z^{\mathrm{der}}_{\mathrm{ch}}(C_X)\) on \(K3\times E\).  The
level-\(\mathsf A\) Hall--Borcherds residual is different: it asks for
the gravity-line operator algebra acting on the boundary and remains
open in Vol~II.

\begin{definition}[Vol~II Hall--Borcherds trace comparison at level \(\mathsf Z\)]
\label{def:G-vol2-hall-borcherds}
\textnormal{[definitional]}
The \emph{level-\(\mathsf Z\) Hall--Borcherds trace comparison} on
\(K3\times E\) is the assertion identifying the protected graded trace
\[
\operatorname{Tr}_{\mathrm{prot}}\bigl((-1)^F
q^{L_0}r^{J_0}s^{\widetilde{L}_0}\bigr)
=Z_{\mathrm{BPS}}^{K3\times E}
=\Delta_5^{-2}
\]
of the retained \(E_3\)-prefactorization algebra
\(\mathcal A_{K3\times E}\) to
the level-\(\mathsf Z\) trace pairing of the chain-level chiral
derived centre \(Z^{\mathrm{der}}_{\mathrm{ch}}(C_X)\) on
\((X=K3\times E,D,\tau)\) at the boundary vacuum \(b\); the
identification \(Z^{\mathrm{der}}_{\mathrm{ch}}(C_X)=
\mathrm{ChirHoch}^\bullet(C_X,C_X)\) is the chiral-Hochschild
concentration theorem of Vol~I (Theorem~H,
\cite{ChiralBarCobarVolI}) on the Koszul locus,
arranged by the level-\(\mathsf Z\) clause~(i) of the cyclic-Hochschild
Beilinson stratification of Vol~II
\cite[chiral-Hochschild Beilinson stratification]{ChiralBarCobarVolII};
the trace pairing of degree \(-3\) is the CY-orientation pairing
clause~(ii) of the same theorem under
\(\mathsf{Eff}_{\mathrm{Pf,or}}\).  The
comparison is a class in
\[
 H^*_{\mathrm{Hall}}(\mathcal A_{K3\times E})
 \;\to\;
 H^*_{\mathrm{Borcherds}}\bigl(Z^{\mathrm{der}}_{\mathrm{ch}}(\mathcal A_{K3\times E})\bigr),
\]
governing whether the chiral-shadow Hall integral
\(\Delta_5^{-2}\) is the shadow of a chain-level chiral-Hochschild
Borcherds integral on the target.
\end{definition}

\begin{remark}[Distinguishing the level-\(\mathsf Z\) trace from the level-\(\mathsf A\) gravity line]
\label{rem:G-level-Z-vs-gravity-line}
The level-\(\mathsf Z\) Hall--Borcherds trace comparison of
Definition~\ref{def:G-vol2-hall-borcherds} has trace-level scope,
whereas level \(\mathsf A\) asks for a
gravity-line operator algebra acting on the boundary chiral algebra of the
\(K3\times E\) holomorphic-topological theory, with Pentagon-face
determinant section \(\Phi_{10}^{\mathrm{un}}=\Delta_5^2\) and
reciprocal partition character \(\Delta_5^{-2}\).  This residual
remains open in Vol~II; its construction requires the operator-valued
completion that the present level-\(\mathsf Z\) trace identity does
not supply.  The level-\(\mathsf Z\) facet is a trace pairing on
\(Z^{\mathrm{der}}_{\mathrm{ch}}(C_X)\); the level-\(\mathsf A\)
facet is an acting algebra on a gravity-line module.  No theorem
constructs that acting algebra.
\end{remark}

\subsection{The level-\texorpdfstring{$\mathsf Z$}{Z} Hall--Borcherds trace comparison on \(K3\times E\)}
\label{app:G-vol2-discharge}

The level-\(\mathsf Z\) Hall--Borcherds trace comparison of
Definition~\ref{def:G-vol2-hall-borcherds} is a criterion relative to
a retained trace datum on \(K3\times E\).  Vol~I Theorem~H
\cite{ChiralBarCobarVolI} and Vol~II's chiral-Hochschild Beilinson
stratification \cite[chiral-Hochschild Beilinson stratification]{ChiralBarCobarVolII}
apply only after that datum has placed the compact source in their
domain.  The level-\(\mathsf A\) operator-algebra facet remains open in
Vol~II (Remark~\ref{rem:G-level-Z-vs-gravity-line}).

\begin{definition}[Retained level-\(\mathsf Z\) Hall--Borcherds trace datum]
\label{def:G-level-Z-trace-datum}
\textnormal{[definitional]}
The datum \((Z_{\mathrm{HB}})\) consists of the following finite and
pro-compatible data:
\begin{enumerate}
\item[\textnormal{(Z1)}] a boundary chart algebra \(A_b\) in the
Vol~I/Vol~II chiral-Koszul domain, together with a filtered
quasi-isomorphism \(A_b\simeq C_X\);
\item[\textnormal{(Z2)}] a class-\(M\) weight completion of
\(\mathrm{Bar}(A_b)\), compatible with the HN inverse system and the
finite Koszul comparison and Mittag--Leffler data;
\item[\textnormal{(Z3)}] the Vol~I chiral-Hochschild concentration
hypothesis on \(A_b\) and the Vol~II stratification hypotheses at
level \(\mathsf Z\);
\item[\textnormal{(Z4)}] a degree \(-3\) cyclic trace pairing on
\(Z^{\mathrm{der}}_{\mathrm{ch}}(A_b)\), compatible with the
Calabi--Yau three-form, the Pfaffian line, and HN transition;
\item[\textnormal{(Z5)}] compatibility of the trace normalization with
the type-II Borcherds product of \(\Delta_5\), the OP chamber scalar
branch, and the protected trace convention of
Chapter~\ref{ch:zbps-realization}.
\end{enumerate}
\((Z_{\mathrm{HB}})\) has trace-level scope:
it supplies a trace pairing on the centre, not an operator algebra
acting on a gravity-line module.
In the present trace packet these rows are not supplied.  The finite
obstruction ledger
\[
\texttt{certificates/trace/k3e\_protected\_trace/blocked\_obligations.csv}
\]
records the missing trace-category, protected-trace-functor,
trace-operator, scalar-normalization, trace-identity, forgetful-map,
gravity-residual, transition, and scalar-firewall rows.  Its verifier
returns \(\texttt{TRACE\_OBSTRUCTION\_LEDGER\_VERIFIED}\), with
\(\texttt{protected\_trace\_certification=false}\) and
\(\texttt{mathematical\_certification=false}\).  Thus the packet is a
finite absence ledger for \((Z_{\mathrm{HB}})\), not a construction of
the level-\(\mathsf Z\) protected trace.
\end{definition}

\begin{theorem}[Level-\(\mathsf Z\) Hall--Borcherds trace criterion on \(K3\times E\)]
\label{thm:G-vol2-discharge}
Granted the retained level-\(\mathsf Z\) trace datum
\((Z_{\mathrm{HB}})\), Vol~I Theorem~H \cite{ChiralBarCobarVolI}, and
Vol~II's chiral-Hochschild Beilinson stratification
\cite{ChiralBarCobarVolII} as established theorems, the
level-\(\mathsf Z\) Hall--Borcherds trace comparison of
Definition~\ref{def:G-vol2-hall-borcherds} holds on \(K3\times E\).
Equivalently, the level-\(n\) protected scalar trace
\(Z_{\mathrm{BPS}}^{K3\times E}=\Delta_5^{-2}\) of
Chapter~\ref{ch:zbps-realization} is identified with the chain-level
CY-orientation trace of the derived chiral centre
\(Z^{\mathrm{der}}_{\mathrm{ch}}(A_b)\) on \(K3\times E\):
\[
\boxed{\;
 \mathrm{Tr}^{\mathrm{bulk}}_n\!
 \bigl((-1)^F\cdot\mathrm{id}\bigr)_{Z^{\mathrm{der}}_{\mathrm{ch}}(C_X)}
 \;=\;
 \Delta_5^{-2}
\;}.
\]
The OP chamber branch is the normalized scalar representative
\(-4096\,\mathrm{Tr}^{\mathrm{bulk}}_n\), not the trace itself.
The comparison is at level \(\mathsf Z\) in the Beilinson chain
of Vol~II; the level-\(\mathsf A\) gravity-line operator algebra is not
constructed here and remains open in Vol~II
(Remark~\ref{rem:G-level-Z-vs-gravity-line}).
\end{theorem}

\begin{proof}
The retained datum \((Z_{\mathrm{HB}})\) supplies the chart algebra
\(A_b\), its comparison with the compact source \(C_X\), the
class-\(M\) completion, the cyclic trace pairing, and the normalization
comparison with the type-II Borcherds product.  Clause~(Z3), together
with Vol~I Theorem~H and the
level-\(\mathsf Z\) clause of Vol~II, identifies the chiral derived
centre with chiral Hochschild cohomology on the retained Koszul locus:
\[
Z^{\mathrm{der}}_{\mathrm{ch}}(A_b)
\;=\;\mathrm{ChirHoch}^\bullet(A_b,A_b),
\]
and clause~(Z4) supplies the degree-\(-3\) trace pairing
\[
\mathrm{Tr}_{\mathcal C}
\colon
Z^{\mathrm{der}}_{\mathrm{ch}}(A_b)\otimes
Z^{\mathrm{der}}_{\mathrm{ch}}(A_b)
\;\longrightarrow\;
\C[-3]
\]
at level \(\mathsf Z\).  Clause~(Z5), the Pfaffian normalization of
Theorem~\ref{thm:ch6-pfaffian-dirac}, and the OP/Igusa scalar
normalization of Theorem~\ref{thm:zbps-level-n-shadow} identify this
trace with the protected scalar:
\[
 \mathrm{Tr}^{\mathrm{bulk}}_n
 =\Delta_5^{-2}.
\]
The OP branch convention of Convention~\ref{conv:op-chamber} gives the
separate scalar representative
\(Z^X_{\mathrm{OP}}=-4096\,\mathrm{Tr}^{\mathrm{bulk}}_n\).
This proves exactly the level-\(\mathsf Z\) criterion: once the retained
trace datum exists with the stated normalization compatibilities, the
trace comparison follows.  The
\(\mathsf{SC}^{\mathrm{ch,top}}\)-brace, the \(E_3\)-topological
factorisation-algebra structure on \(X\times\mathbb R\), and the
gravity-line operator algebra are not consequences of this trace datum.
\end{proof}

\begin{theorem}[Pfaffian--Dirac theorem and level-\(\mathsf Z\) Z\(_{\mathrm{BPS}}\) trace identity on \(K3\times E\), relative to retained data]
\label{thm:G-pfaffian-dirac-relative}
Granted Theorems~\ref{thm:G-D0}, \ref{thm:G-O1},
\ref{thm:G-O1-plus}, \ref{thm:G-O2-orbit-transport}, the retained
finite geometric Pfaffian datum \((P_{\mathrm{fin}})\), and
\ref{thm:G-vol2-discharge}, the last itself granted
\((Z_{\mathrm{HB}})\), Vol~I
Theorem~H \cite{ChiralBarCobarVolI} and Vol~II's chiral-Hochschild
Beilinson stratification \cite{ChiralBarCobarVolII} as established
theorems, the Pfaffian--Dirac theorem~\ref{thm:ch6-pfaffian-dirac},
the orientation-character match~\ref{thm:ch6-orientation-character},
and the level-\(\mathsf Z\) trace identity for the protected scalar
\(Z_{\mathrm{BPS}}^{K3\times E}=\Delta_5^{-2}\) to a chain-level
CY-orientation trace on \(Z^{\mathrm{der}}_{\mathrm{ch}}(C_X)\) all
hold on \(K3\times E\).  The primitive recognition theorem
\ref{thm:ch6-primitive-recognition} is a separate conditional theorem:
it requires the compact-source representatives, relations, pairings,
radicals, primitive-space and pairing-kernel Mittag--Leffler exactness,
PBW comparison with PBW associated-graded Mittag--Leffler exactness, and
full parity dimensions of \((S1)\)--\((S10)\).
The level-\(\mathsf A\) gravity-line operator algebra
remains open in Vol~II
(Remark~\ref{rem:G-level-Z-vs-gravity-line}).
\end{theorem}

\begin{proof}
\((D0)\) follows from the retained D0-HN datum through
Theorem~\ref{thm:G-D0};
\((O1)\) follows from the retained quotient-orientation datum through
Theorem~\ref{thm:G-O1}; the equivariance part of \((O1)^+\) follows
from the chosen Weyl lifts through
Theorem~\ref{thm:G-O1-plus}; and \((O2)\) is transported from the
retained wall-atlas datum by Theorem~\ref{thm:G-O2-orbit-transport}.
The finite geometric Pfaffian data are the datum \((P_{\mathrm{fin}})\)
of Definition~\ref{def:G-finite-pfaffian-section}.  The
level-\(\mathsf Z\) Hall--Borcherds trace comparison follows from the
retained \((Z_{\mathrm{HB}})\) datum through
Theorem~\ref{thm:G-vol2-discharge}.  The Pfaffian and orientation
theorems of Chapter~\ref{ch:pfaffian-dirac} then hold on
\(K3\times E\) relative to these retained data.  The
primitive-recognition theorem asks for the
additional compact-source datum \((S1)\)--\((S10)\), and Appendix~G
gives the Pfaffian-orientation lane, not the compact-source
Lie-superalgebra recognition lane.  The level-\(\mathsf Z\) trace
identity of Chapter~\ref{ch:zbps-realization} reads
\(\mathrm{Tr}^{\mathrm{bulk}}_n=\Delta_5^{-2}\).  The OP scalar branch
is \(-4096\,\mathrm{Tr}^{\mathrm{bulk}}_n\), by
Convention~\ref{conv:op-chamber}.
\end{proof}

\begin{remark}[The conditioning is the retained level-\(\mathsf Z\) datum and Vol~I/Vol~II chain-level theorems]
\label{rem:G-conditioning-clean}
After the \(K3\times E\) retained D0-HN datum, the retained
\((O1_{\mathrm{quot}})\) datum, the chosen Weyl lifts, the retained
\((O2_{\mathrm{atlas}})\) datum, and the trace comparison of
\S\ref{app:G-vol2-residual}, the remaining chain-level conditioning
for the level-\(\mathsf Z\) Z\(_{\mathrm{BPS}}\) trace identity is the
retained datum \((Z_{\mathrm{HB}})\), Vol~I Theorem~H
\cite{ChiralBarCobarVolI} (chiral Hochschild concentration on the
Koszul locus), and Vol~II
\cite[chiral-Hochschild Beilinson stratification]{ChiralBarCobarVolII}
clauses~(i) and~(ii) (chain-level bulk identification and
CY-orientation trace pairing) as correct theorems.  The
K3\(\times E\)-specific trace obstruction is precisely
\((Z_{\mathrm{HB}})\).  The theorem is conditional on this retained
datum and on the Vol~I/Vol~II chain-level constructions.
The finite packet
\(\texttt{certificates/trace/k3e\_protected\_trace}\) presently verifies
only the absence ledger
\(\texttt{TRACE\_OBSTRUCTION\_LEDGER\_VERIFIED}\), with
\(\texttt{protected\_trace\_certification=false}\); it does not supply
\((Z_{\mathrm{HB}})\).
The level-\(\mathsf A\) gravity-line operator algebra sits outside this
conditioning and remains open in Vol~II.
Primitive recognition is independent of the retained orientation data:
it is exactly the
finite compact-source recognition problem of
Theorem~\ref{thm:ch6-primitive-recognition}.
\end{remark}

\subsection{Compatibility of the retained data}
\label{app:G-retained-data-compatibility}

The retained D0-HN datum, the retained quotient-orientation datum, the
chosen Weyl lifts, and the retained \((O2_{\mathrm{atlas}})\) datum convert
the Pfaffian and orientation parts of the Pfaffian--Dirac theorem to
statements on \(K3\times E\) at the Pfaffian and orientation level.
The primitive
Lie-superalgebra recognition remains conditional on \((S1)\)--\((S10)\).
The proof uses the retained data in a fixed implication chain.
The retained D0-HN datum gives Theorem~\ref{thm:G-D0}.  The retained
quotient-orientation datum gives Theorem~\ref{thm:G-O1}.  The Weyl
lifts and simple-wall signs give Theorem~\ref{thm:G-O1-plus}.  The
retained \((O2_{\mathrm{atlas}})\) datum together with these Weyl lifts
gives Theorem~\ref{thm:G-O2-orbit-transport}.  The scalar separation
\((R1)\)--\((R2)\) is Theorem~\ref{thm:G-R1-R2}.  The
Hall--Borcherds trace comparison is Theorem~\ref{thm:G-vol2-discharge}.
The gravity-line algebra is a separate Vol~II construction.

The criteria use the following theorems, with the fourfold analogue
separated from the \(K3\times E\) argument:
\begin{itemize}
\item Joyce--Upmeier orientation theorem
\cite[Theorem~1.11]{JoyceUpmeier};
\item Joyce--Tanaka--Upmeier quasi-projective extension
\cite[Theorem~1.4]{JoyceTanakaUpmeier};
\item the Cao--Kool--Monavari CY\(_4\) cosection-twist extension
\cite[Theorem~1.6]{CaoKoolMonavari} is a fourfold analogue, not a
theorem used in the \(K3\times E\) CY\(_3\) reduction above;
\item Kiem--Li cosection localisation
\cite[Theorem~5.0.1]{KiemLi2013};
\item Brav--Bussi--Dupont--Joyce--Szendr\H{o}i d-critical
chart cocycle \cite[Theorem~6.9]{BBDJS2015};
\item Pantev--To\"en--Vaqui\'e--Vezzosi
\((-1)\)-shifted symplectic structure
\cite[Theorem~0.3]{PTVV2013};
\item Maulik--Pandharipande K3 reduced theory
\cite[Theorem~1]{MaulikPandharipande};
\item Oberdieck--Pandharipande \(K3\times E\) Igusa cusp form
factorisation \cite[Theorem~0.2]{OberdieckK3xE},
\cite[Theorem~1]{OB}.
\end{itemize}

The \(K3\times E\)-specific finite-stabilizer criteria are
Lemma~\ref{lem:G-Klein-four} (Klein-four detection on
rank-\(2\) E-strata), Lemma~\ref{lem:G-two-primary-stabilizer}
(full two-primary \(E[N]\)-edge detection on finite-stabilizer strata), and
Theorem~\ref{thm:G-R1-R2} (\((R1)\)--\((R2)\)
scalar separation).  The foundational citations give
the chart and orientation technology; Mittag--Leffler exactness is a
condition in the retained datum \((D0_{\mathrm{HN}})\).

\subsection{The relative theorem on K3$\times$E}
\label{app:G-relative-theorem}

The Pfaffian theorem~\ref{thm:ch6-pfaffian-dirac} and the
orientation-character theorem~\ref{thm:ch6-orientation-character} hold
on \(K3\times E\) at the Pfaffian and orientation level relative to
\((D0_{\mathrm{HN}})\), \((O1_{\mathrm{quot}})\), the Weyl lifts,
\((O2_{\mathrm{atlas}})\), and \((P_{\mathrm{fin}})\).  The obstruction \((D0)\) of
Chapter~\ref{chap:dirac-igusa-pro-object} is reduced to the retained
D0-HN datum by Theorem~\ref{thm:G-D0}; \((O1)\) and \((O1)^+\) are reduced through
Theorems~\ref{thm:G-O1} and \ref{thm:G-O1-plus} relative to their
retained data; \((O2)\) is transported from \((O2_{\mathrm{atlas}})\)
by Theorem~\ref{thm:G-O2-orbit-transport}.  The
\((R1)\)--\((R2)\) scalar separation
\(4096=2^{12}=64^2\) is
Theorem~\ref{thm:G-R1-R2}.  The level-\(\mathsf Z\) Hall--Borcherds
trace identity
\(\mathrm{Tr}^{\mathrm{bulk}}_n=\Delta_5^{-2}\) on
\(Z^{\mathrm{der}}_{\mathrm{ch}}(C_X)\) is
Theorem~\ref{thm:G-vol2-discharge}, granted
\((Z_{\mathrm{HB}})\), Vol~I Theorem~H
\cite{ChiralBarCobarVolI}, and Vol~II's chiral-Hochschild Beilinson
stratification \cite{ChiralBarCobarVolII} as established theorems;
no appeal is made to the Vol~II Universal Holography master theorem,
whose standard-landscape scope (affine Kac--Moody at non-critical
level, principal \(W_{N,c}\), \(\mathrm{Vir}_c\), Schellekens
\(c=24\), Monster \(V^\natural\), VSKR+BGG-tempered irrational
cosets on a smooth projective curve) does not include the
BKM-derived hybrid factorisation chart \(C_X\) on
\(K3\times E\).  The Saito--Kurokawa eigenline arithmetic of
Proposition~\ref{prop:zbps-sk-hecke} and the \(\kappa\)-universality
identity \(\kappa_{\mathrm{BKM}}(\Phi_N)=c_N(0)/2\) across the CHL
frame \(N\in\{1,2,3,4,6\}\) of
Theorem~\ref{thm:bkm-kappa-universal} are inscribed by direct
computation.  Throughout, the cited foundational results
(Joyce--Upmeier, BBDJS,
Maulik--Pandharipande, Oberdieck--Pandharipande, Kiem--Li,
Pantev--To\"en--Vaqui\'e--Vezzosi) are applied in their published
form.  The Cao--Kool--Monavari CY\(_4\) theorem is a separate
fourfold analogue.

The theorem-level automorphic, BKM-denominator, and singular-theta
readings of \(\Delta_5\) close at the automorphic level, and the
protected Pfaffian
\(\Pf_{\mathrm{prot}}(\mathfrak D_X^{\mathrm{DI}})\) is identified with
\(\Delta_5\) relative to the retained data.  The Lie-superalgebra
recognition \(P_X^{\Pi,\mathrm{red}}\cong\mathfrak g_{\Delta_5}\)
requires the finite compact-source datum \((S1)\)--\((S10)\).  The
mirror discriminant of Chapter~\ref{ch:view4-mirror} is the remaining
view: the automorphic calculation gives
\([\Delta_5^{-2}]=2[\mathfrak H_{\mathrm{II}}]\), while the
identification of this Heegner divisor with the reduced mirror discriminant and
the global trivialisation of the comparison line bundle are conjectural
and refined into four period conditions in
Conjecture~\ref{conj:mirror-discriminant}.

The unresolved data are the finite compact-source recognition datum
\((S1)\)--\((S10)\), the mirror period comparison of
Conjecture~\ref{conj:mirror-discriminant}, the explicit OPE structure
on \(Z^{\mathrm{der}}_{\mathrm{ch}}(C_X)\), the Type~II sigma-model on
\(K3\times E\) with worldsheet BPS partition function \(\Delta_5^{-2}\),
the level-\(\mathsf A\) gravity-line operator algebra, and the
Mathieu--umbral twined denominator system.  For the order-seven and
order-fifteen conjugacy classes \(\{7A,7B,15A,15B\}\) of \(M_{24}\),
Conjecture~\ref{conj:order-seven-fifteen-twined-denominators} of
Chapter~\ref{ch:open-obstructions} refines the
Cheng--Duncan--Harvey--Gannon umbral moonshine conjecture to four
assertions on
\(\Theta(\phi^{(g),\mathrm{Weil}}_{0,1})\!\bigm|_{\Hpl_2}\), and
identifies the arithmetic obstruction in the parity of \(c_g(0)\).
The level-\(\mathsf A\) gravity-line problem demands an operator algebra
acting on the boundary chiral algebra of the \(K3\times E\)
holomorphic-topological theory with Pentagon-face determinant section
\(\Phi_{10}^{\mathrm{un}}=\Delta_5^2\) and reciprocal partition
character \(\Delta_5^{-2}\), an operator-valued
Hall--Borcherds completion, and the 6d holomorphic
Chern--Simons quartic anomaly vanishing on \(K3\times E\) of Vol~III.
The level-\(\mathsf Z\) trace criterion of
Theorem~\ref{thm:G-vol2-discharge} is a trace pairing on the centre,
not an acting algebra on a gravity-line module.
